\section{RESPUESTAS}
\subsection{Ejercicio 2.1}
\subsubsection{Tabla de verdad}
\sangria{} Un BCD transforma los dígitos en notación decimal (del 0 al 9), en binario de 4 bits. Entonces:
\begin{center}
\begin{tabular}{c|cccc}
\rowcolor[HTML]{FFFC9E} 
N. en decimal & A & B & C & D \\ \hline
0             & 0 & 0 & 0 & 0 \\
1             & 0 & 0 & 0 & 1 \\
2             & 0 & 0 & 1 & 0 \\
3             & 0 & 0 & 1 & 1 \\
4             & 0 & 1 & 0 & 0 \\
5             & 0 & 1 & 0 & 1 \\
6             & 0 & 1 & 1 & 0 \\
7             & 0 & 1 & 1 & 1 \\
8             & 1 & 0 & 0 & 0 \\
9             & 1 & 0 & 0 & 1
\end{tabular}
\end{center}
Siendo $A,B,C,D$ las posiciones de los 4 bits. EL Exceso-3 es un código que adiciona $3$ a un valor de entrada en decimal:
\ecuacion{XS3 = Decimal + 3}
En binario, el Exceso-3 se ve de la misma forma:
\begin{center}
$0000 \rightarrow 0011$ \\
$0001 \rightarrow 0100$ \\
$0010 \rightarrow 0101$ \\
$...$
\end{center}
\noindent{}Hemos definido 4 funciones, cada una representa los bits de Exceso-3:
\begin{center}
    $W(A,B,C,D) \quad X(A,B,C,D) \quad Y(A,B,C,D) \quad Z(A,B,C,D)$
\end{center}

\noindent{} Para cada combinación de los bits $A,B,C\text{ y }D$, se armará otra combinación de 4 bits en código Exceso-3:
\begin{center}
\begin{tabular}{c|cccc|l|cccc|c}
\cellcolor[HTML]{FFFC9E}N. en decimal & \cellcolor[HTML]{FFFC9E}A & \cellcolor[HTML]{FFFC9E}B & \cellcolor[HTML]{FFFC9E}C & \cellcolor[HTML]{FFFC9E}D &               & \cellcolor[HTML]{FFCCC9}W & \cellcolor[HTML]{FFCCC9}X & \cellcolor[HTML]{FFCCC9}Y & \cellcolor[HTML]{FFCCC9}Z & \cellcolor[HTML]{FFCCC9}N. en decimal \\ \cline{1-5} \cline{7-11} 
0                                     & 0                         & 0                         & 0                         & 0                         & $\rightarrow$ & 0                         & 0                         & 1                         & 1                         & 3                                     \\
1                                     & 0                         & 0                         & 0                         & 1                         & $\rightarrow$ & 0                         & 1                         & 0                         & 0                         & 4                                     \\
2                                     & 0                         & 0                         & 1                         & 0                         & $\rightarrow$ & 0                         & 1                         & 0                         & 1                         & 5                                    
\end{tabular}
\end{center}

\noindent La misma regla continua para el resto de combinaciones de bits para el $BCD$, quedando conformada la tabla de verdad solicitada:
\begin{center}
\begin{tabular}{c|cccc|cccc|c}
\rowcolor[HTML]{FFFC9E} 
N. en decimal & A & B & C & D & \cellcolor[HTML]{FFCCC9}W(A,B,C,D) & \cellcolor[HTML]{FFCCC9}X(A,B,C,D) & \cellcolor[HTML]{FFCCC9}Y(A,B,C,D) & \cellcolor[HTML]{FFCCC9}Z(A,B,C,D) & \cellcolor[HTML]{FFCE93}Decimal exceso-3 \\ \hline
0             & 0 & 0 & 0 & 0 & 0                                  & 0                                  & 1                                  & 1                                  & 3                                        \\
1             & 0 & 0 & 0 & 1 & 0                                  & 1                                  & 0                                  & 0                                  & 4                                        \\
2             & 0 & 0 & 1 & 0 & 0                                  & 1                                  & 0                                  & 1                                  & 5                                        \\
3             & 0 & 0 & 1 & 1 & 0                                  & 1                                  & 1                                  & 0                                  & 6                                        \\
4             & 0 & 1 & 0 & 0 & 0                                  & 1                                  & 1                                  & 1                                  & 7                                        \\
5             & 0 & 1 & 0 & 1 & 1                                  & 0                                  & 0                                  & 0                                  & 8                                        \\
6             & 0 & 1 & 1 & 0 & 1                                  & 0                                  & 0                                  & 1                                  & 9                                        \\
7             & 0 & 1 & 1 & 1 & 1                                  & 0                                  & 1                                  & 0                                  & 10                                       \\
8             & 1 & 0 & 0 & 0 & 1                                  & 0                                  & 1                                  & 1                                  & 11                                       \\
9             & 1 & 0 & 0 & 1 & 1                                  & 1                                  & 0                                  & 0                                  & 12                                      
\end{tabular}
\end{center}

\subsubsection{Funciones canónicas}
\sangria{} La función canónica es aquella donde aparecen todos los términos de los cuales ella depende, en formato de minitérminos o maxitérminos. En este caso se usó minitérminos (suma de productos, tomamos los 1 de la función).

\begin{center}
\begin{minipage}{0.45\textwidth}\centering
\begin{tabular}{cccc|c}
\rowcolor[HTML]{FFFC9E} 
A & B & C & D & \cellcolor[HTML]{FFCCC9}W(A,B,C,D) \\ \hline
0 & 0 & 0 & 0 & 0                                  \\
0 & 0 & 0 & 1 & 0                                  \\
0 & 0 & 1 & 0 & 0                                  \\
0 & 0 & 1 & 1 & 0                                  \\
0 & 1 & 0 & 0 & 0                                  \\
0 & 1 & 0 & 1 & 1                                  \\
0 & 1 & 1 & 0 & 1                                  \\
0 & 1 & 1 & 1 & 1                                  \\
1 & 0 & 0 & 0 & 1                                  \\
1 & 0 & 0 & 1 & 1                                 
\end{tabular}\\[10pt]
$W = \overline{A}\cdot{}B\cdot{}\overline{C}\cdot{}D + \overline{A}\cdot{}B\cdot{}C\cdot\overline{D} + \overline{A}\cdot{}B\cdot{}C\cdot{}D + A\cdot\overline{B}\cdot\overline{C}\cdot\overline{D} + A\cdot\overline{B}\cdot\overline{C}\cdot{}D$
\end{minipage}
\hfill
\begin{minipage}{0.45\textwidth}\centering
\begin{tabular}{cccc|c}
\rowcolor[HTML]{FFFC9E} 
A & B & C & D & \cellcolor[HTML]{FFCCC9}X(A,B,C,D) \\ \hline
0 & 0 & 0 & 0 & 0                                  \\
0 & 0 & 0 & 1 & 1                                  \\
0 & 0 & 1 & 0 & 1                                  \\
0 & 0 & 1 & 1 & 1                                  \\
0 & 1 & 0 & 0 & 1                                  \\
0 & 1 & 0 & 1 & 0                                  \\
0 & 1 & 1 & 0 & 0                                  \\
0 & 1 & 1 & 1 & 0                                  \\
1 & 0 & 0 & 0 & 0                                  \\
1 & 0 & 0 & 1 & 1                                 
\end{tabular}\\[10pt]
$X = \overline{A}\cdot\overline{B}\cdot\overline{C}\cdot{}D + \overline{A}\cdot\overline{B}\cdot{}C\overline{D} + \overline{A}\cdot\overline{B}\cdot{}C\cdot{}D + \overline{A}\cdot{}B\cdot\overline{C}\cdot\overline{D} + A\cdot\overline{B}\cdot\overline{C}\cdot{}D$
\end{minipage}
\end{center}
\newpage
\begin{center}
\begin{minipage}{0.45\textwidth}\centering
\begin{tabular}{cccc|c}
\rowcolor[HTML]{FFFC9E} 
A & B & C & D & \cellcolor[HTML]{FFCCC9}Y(A,B,C,D) \\ \hline
0 & 0 & 0 & 0 & 1                                  \\
0 & 0 & 0 & 1 & 0                                  \\
0 & 0 & 1 & 0 & 0                                  \\
0 & 0 & 1 & 1 & 1                                  \\
0 & 1 & 0 & 0 & 1                                  \\
0 & 1 & 0 & 1 & 0                                  \\
0 & 1 & 1 & 0 & 0                                  \\
0 & 1 & 1 & 1 & 1                                  \\
1 & 0 & 0 & 0 & 1                                  \\
1 & 0 & 0 & 1 & 0                                 
\end{tabular}\\[10pt]

$Y = \overline{A}\cdot\overline{B}\cdot\overline{C}\cdot\overline{D} + \overline{A}\cdot\overline{B}\cdot{}C\cdot{}D + \overline{A}\cdot{}B\cdot\overline{C}\cdot\overline{D} + \overline{A}\cdot{}B\cdot{}C\cdot{}D + A\cdot\overline{B}\cdot\overline{C}\cdot\overline{D}$
\end{minipage}
\begin{minipage}{0.45\textwidth}\centering
\begin{tabular}{cccc|c}
\rowcolor[HTML]{FFFC9E} 
A & B & C & D & \cellcolor[HTML]{FFCCC9}Z(A,B,C,D) \\ \hline
0 & 0 & 0 & 0 & 1                                  \\
0 & 0 & 0 & 1 & 0                                  \\
0 & 0 & 1 & 0 & 1                                  \\
0 & 0 & 1 & 1 & 0                                  \\
0 & 1 & 0 & 0 & 1                                  \\
0 & 1 & 0 & 1 & 0                                  \\
0 & 1 & 1 & 0 & 1                                  \\
0 & 1 & 1 & 1 & 0                                  \\
1 & 0 & 0 & 0 & 1                                  \\
1 & 0 & 0 & 1 & 0                                 
\end{tabular} \\[10pt] $Z = \overline{A}\cdot\overline{B}\cdot\overline{C}\cdot\overline{D} + \overline{A}\cdot\overline{B}\cdot{}C\cdot\overline{D} + \overline{A}\cdot{}B\cdot\overline{C}\cdot\overline{D} + \overline{A}\cdot{}B\cdot{}C\cdot\overline{D} + A\cdot\overline{B}\cdot\overline{C}\cdot\overline{D}$
\end{minipage}
\end{center}
\subsubsection{Minimización}
\sangria{} Las minimizaciones permite quitar redundancia de compuertas en la función, y el resultado es una versión con menos compuertas de la misma. Para ello se ha usado el mapa de Karnaugh:

\begin{center}\begin{minipage}{0.45\textwidth}\centering\begin{karnaugh-map}[4][4][1][$C\qquad{}D$][$A\qquad{}B$]
    \minterms{5,6,7,8,9}
    \maxterms{0,1,2,3,4}
    \terms{12,13,15,14,11,10}{$x$}
    \implicant{0}{2}\implicant{0}{4}
    \end{karnaugh-map}$W = (A + B)\cdot{}(A + C +D)$\end{minipage}
\begin{minipage}{0.45\textwidth}\centering
    \begin{karnaugh-map}[4][4][1][$C\qquad{}D$][$A\qquad{}B$]
    \minterms{1,2,3,4,9}\maxterms{0,5,6,7,8}
    \terms{12,13,15,14,11,10}{$x$}
    \implicant{5}{15}\implicant{7}{14}
    \implicantedge{0}{0}{8}{8}
    \end{karnaugh-map}
    $X = (B + C +D )\cdot{}(\overline{B} + \overline{D})\cdot{}(\overline{B} + \overline{C})$
\end{minipage}\end{center}

\begin{center}\begin{minipage}{0.45\textwidth}\centering\begin{karnaugh-map}[4][4][1][$C\qquad{}D$][$A\qquad{}B$]
    \minterms{0,3,4,7,8}
    \maxterms{1,2,5,6,9}
    \terms{12,13,15,14,11,10}{$x$}
    \implicant{0}{8}\implicant{3}{11}
\end{karnaugh-map}
$Y = \overline{C}\cdot{}\overline{D} + C\cdot{}D$
\end{minipage}\begin{minipage}{0.45\textwidth}\centering\begin{karnaugh-map}[4][4][1][$C\qquad{}D$][$A\qquad{}B$]
    \minterms{0,2,4,6,8}
    \maxterms{1,3,5,7,9}
    \terms{12,13,15,14,11,10}{$x$}
    \implicant{0}{8}\implicant{2}{10}
\end{karnaugh-map}
$Z = \overline{C}\cdot\overline{D} + \overline{C}\cdot{}D$
\end{minipage}
\end{center}


